\begin{pdSolution}{bonded:ex:bonded-cap-attenuation}
The rejected card violates \emph{temporal} attenuation: Theorem~\ref{bonded:thm:scope} requires $C(\tau_{k}) \subseteq C(\tau_{k-1})$ along every dimension of the capability, and TTL is one such dimension---a child's expiry must not exceed its parent's. Here the child requests $20$~minutes against a parent that expires at $15$; even though the capability set $\{\texttt{fs:write:auth.ts}\}$ is a subset of $\alpha$'s, the delegation as a whole is not a narrowing.

An acceptable card is $\{\texttt{fs:write:auth.ts}\}$, TTL~$\le 15$~minutes---dropping \texttt{db:read} (a valid narrowing of the capability set) and matching or shortening the TTL. A strictly narrower example is $\{\texttt{fs:read:auth.ts}\}$, TTL~$=10$~minutes, if the deployment's capability lattice places \texttt{read} below \texttt{write} on the same resource.
\end{pdSolution}
\begin{pdSolution}{bonded:ex:bonded-advisory-claims}
An advisory claim is a durable declaration that an agent intends to use a resource; it informs coordination (Definition~\ref{bonded:def:advisory}'s conflict graph) but does not make conflicting effects impossible. Every trace an enforced-lock regime allows, an advisory regime also allows---the lock regime is a special case that additionally forbids the conflicting trace---so advisory claims are strictly more permissive, never less.

That extra permissiveness is the feature Theorem~\ref{bonded:prop:sen} names: Sen's theorem shows that granting an agent a decisive personal domain (an enforced lock) can block a Pareto-superior team outcome the allocator cannot see coming. An advisory regime never grants that veto, so it never forecloses the outcome a lock would have blocked---at the cost of also admitting traces where two agents genuinely do collide. \S\ref{bonded:sec:governance}'s appeals and preemption path is exactly the mechanism that prices the cost of that added permissiveness rather than eliminating it.
\end{pdSolution}
\begin{pdSolution}{bonded:ex:bonded-reputation-vs-bonds}
First, reputation compensates nobody for the harm already done. A one-shot defector takes the gain from the first (and only) defection and leaves before any future exclusion can matter to it---exactly the ``one-shot defectors'' adversary class named above. Second, cheap identities let the actor whitewash the loss: register a fresh agent ID, and the reputation ledger has nothing to slash. A pre-funded bond avoids both failure modes because it creates recoverable value \emph{before} authority is granted and keys the exposure to the durable principal (\S\ref{bonded:sec:attribution-survives}) rather than to the disposable agent identity---it prices and contains harm at the moment of the first defection, not on the promise of a future one. What a bond does \emph{not} do is prove the work will be correct; it only ensures that if the work is wrong, someone has already paid for the possibility.
\end{pdSolution}
\begin{pdSolution}{bonded:ex:bonded-claim-signaling-pd}
Fix $A$'s choice and compare its two rows against each of $B$'s columns. Against $B{:}\mathrm{T}$, $A$ gets $3$ from $\mathrm{T}$ and $4$ from $\mathrm{F}$, so $\mathrm{F}$ is better ($4>3$). Against $B{:}\mathrm{F}$, $A$ gets $0$ from $\mathrm{T}$ and $1$ from $\mathrm{F}$, so $\mathrm{F}$ is again better ($1>0$). $\mathrm{F}$ therefore strictly dominates $\mathrm{T}$ for $A$ regardless of $B$'s move; the game is symmetric, so the same holds for $B$. Checking the three remaining cells: at $(T,T)=(3,3)$ either player can deviate to $\mathrm{F}$ and gain ($4>3$); at $(T,F)=(0,4)$, $A$ can deviate to $\mathrm{F}$ and gain ($1>0$); at $(F,T)=(4,0)$, $B$ can deviate to $\mathrm{F}$ and gain ($1>0$). Only $(F,F)=(1,1)$ survives---neither player gains by unilaterally switching back to $\mathrm{T}$ ($0<1$)---so it is the unique Nash equilibrium, exactly as the table's caption states.

\emph{Provenance note.} An earlier draft of this chapter's stage-game table used different payoffs that a review flagged as \emph{not} a genuine Prisoner's Dilemma---mutual truth was not Pareto-dominant, and the proposed mutual-punishment strategy was not sequentially rational under that table. Table~\ref{bonded:tab:claim-stage-game} above is the corrected table; this exercise verifies the correction that is already in the chapter, not the flawed table that review replaced.
\end{pdSolution}
\begin{pdSolution}{bonded:ex:bonded-threshold-hand}
At $\delta=0.34$: $\delta^2=0.1156$, $1+\delta+\delta^2=1.4556$, $2\delta=0.68$, and $0.68\times1.4556=0.989808$, so $f(0.34)=0.989808-1=-0.010192$ [verified]---negative. At $\delta=0.35$: $\delta^2=0.1225$, $1+\delta+\delta^2=1.4725$, $2\delta=0.7$, and $0.7\times1.4725=1.03075$, so $f(0.35)=1.03075-1=0.03075$ [verified]---positive. $f$ is continuous and strictly increasing on $(0,1)$ (every term's coefficient is positive), so a sign change from $\delta=0.34$ to $\delta=0.35$ proves a root lies strictly between them---the Intermediate Value Theorem, not a numerical coincidence.

This matches two independent committed artifacts. \path{proofs/economics/delta-threshold.z3}'s second query (\path{delta-threshold.expected.txt}) asserts the root lies in $[0.34, 0.35]$ and returns \texttt{sat}, agreeing with the sign change computed here. Independently, \path{proofs/economics/sweep-delta.sh}'s committed sample sweep (documented in \path{proofs/economics/README.md}) reports the model-checked IC invariant \texttt{VIOLATED} at $\delta=0.34$ and \texttt{HOLDS} at $\delta=0.35$---the same sign change, checked by a state machine instead of a closed-form inequality.
\end{pdSolution}
\begin{pdSolution}{bonded:ex:bonded-z3-verdicts}
The script asks, in order: (1) does a real root of $2\delta^3+2\delta^2+2\delta-1=0$ exist in $[0,1]$; (2) does that root lie in the tighter interval $[0.34, 0.35]$; (3) does a \emph{second}, distinct root exist in $[0,1]$ (a uniqueness check, expected to fail).

Pinned verdict sequence (\path{proofs/economics/delta-threshold.expected.txt}):
\begin{lstlisting}
sat
(
  (define-fun delta () Real
    (root-obj (+ (* 2 (^ x 3)) (* 2 (^ x 2)) (* 2 x) (- 1)) 1))
)
sat
(
  (define-fun delta () Real
    (root-obj (+ (* 2 (^ x 3)) (* 2 (^ x 2)) (* 2 x) (- 1)) 1))
)
unsat
\end{lstlisting}
Reading the three lines against the three questions: \texttt{sat} (a root exists, printed in Z3's algebraic \texttt{root-obj} form rather than a decimal); \texttt{sat} again with the identical model (that same root also satisfies the tighter bound, so it lies in $[0.34,0.35]$); \texttt{unsat} (no second, distinct $\delta_2$ in $[0,1]$ satisfies the cubic, so the root is unique). All three together are exactly Theorem~\ref{bonded:prop:claim-signaling-ic}'s threshold, existence-and-uniqueness, mechanized.
\end{pdSolution}
\begin{pdSolution}{bonded:ex:bonded-tla-two-deltas}
At $\delta=0.30$: $\delta^2=0.09$, $1+\delta+\delta^2=1.39$, $2\delta=0.6$, and $0.6\times1.39=0.834$ [verified]. The inequality $1<0.834$ is false, so the deviation is profitable and the invariant \texttt{NoUnilateralDeviationPositive} fails---matching the table's stated verdict and the committed sweep row \texttt{0.30 VIOLATED} (\texttt{proofs/economics/README.md}). At $\delta=0.35$ (\texttt{proofs/economics/claim\_signaling.cfg}'s own default, \texttt{DeltaNum=35}, \texttt{DeltaDen=100}), $2\delta(1+\delta+\delta^2)=1.03075$ from the worked example above, so $1<1.03075$ holds and the invariant holds---matching \texttt{claim\_signaling.cfg} being the one committed configuration and the sweep row \texttt{0.35 HOLDS}.

The trace: at $\delta=0.30$, TLC's search reaches a state where an agent deviated to $F$ once in round one, was punished for three rounds at the mutual-claim payoff, and reached \texttt{round = Horizon} with \texttt{deviated} true; because the deviator's discounted score from that trajectory (\texttt{actualScore}) exceeds the score it would have earned by cooperating throughout (\texttt{followScore}), the state violates \texttt{No}\allowbreak\texttt{Unilateral}\allowbreak\texttt{Deviation}\allowbreak\texttt{Positive} and TLC reports it as a counterexample. No such state exists at $\delta=0.35$: the same one-shot-deviation trajectory now scores strictly worse than cooperating, so every reachable state at \texttt{round=Horizon} satisfies \texttt{actualScore} $\le$ \texttt{followScore} and the search completes with ``No error has been found.'' The committed trace is \path{proofs/economics/claim_signaling_delta30.run.log}, produced by the negative-control configuration \path{claim_signaling_delta30.cfg}; the proofs workflow asserts on every run that it still violates.
\end{pdSolution}
\begin{pdSolution}{bonded:ex:bonded-sweep-crossover}
The committed sample output (\texttt{proofs/economics/README.md}) reports the invariant \texttt{VIOLATED} for every $\delta$ from $0.30$ through $0.34$ and \texttt{HOLDS} from $0.35$ through $0.40$, so the crossover---the smallest $\delta$ for which the invariant holds---is pinned at
\begin{lstlisting}
crossover (smallest delta where invariant HOLDS) = 0.35
closed-form root (delta-threshold.z3)            = 0.3425
PASS: crossover matches closed-form within integer-grid rounding.
\end{lstlisting}
The script's own pass criterion (its exit code) is that the crossover lands in $\{0.34, 0.35\}$---the two integer-grid points bracketing the closed-form root $\delta^\star\approx0.3425$ from Exercise~\ref{bonded:ex:bonded-threshold-hand}. A crossover anywhere else would mean the model-checked state machine and the closed-form cubic disagree about where cooperation becomes sustainable, which is exactly the regression this script exists to catch on every PR touching \texttt{proofs/economics/} (\S\ref{bonded:app:mech}).
\end{pdSolution}
\begin{pdSolution}{bonded:ex:bonded-a5-coverage-bound}
The attacker spawns many low-history Sybil insurer identities, each posting the mandated deposit $B_{\mathrm{dep}}$; each identity underbids honest insurers by an arbitrarily small $\varepsilon$, wins the auction, collects premiums on every no-loss round, and defaults the moment a real loss lands. Because the slash is capped at $\min(B_{\mathrm{dep}}, B_T)$, once $B_{\mathrm{dep}} \ge B_T$ every dollar of deposit beyond the coverage ceiling is recovered intact on no-loss transactions---it never reaches the commons. Concretely, if concurrent aggregate exposure across the Sybil's identities is not reserved, one deposit can underwrite several simultaneous liabilities at once, multiplying the attack for free.

Correction to the naive fix: full reimbursement of a single loss does not, by itself, show the attacker profits---that also needs the premium/loss distribution, exposure reuse across identities, or an external benefit from sabotage. Raising $B_{\mathrm{dep}}$ alone does not close A5, because the slash cap still bounds the attacker's downside at $B_T$ regardless of how large a deposit sits above it. The defense that actually closes it is the composite mechanism named in the hint: an unrecoverable per-identity onboarding cost $C_{\mathrm{kyc}}$ (so spawning $K$ Sybils costs $K \cdot C_{\mathrm{kyc}}$, not zero), reputation gating that caps a new identity to low-coverage transactions until it accumulates a track record, a per-risk-class $B_{\mathrm{dep}}$ tuned to saturate its own coverage cap exactly, and atomic reservation so aggregate outstanding exposure across an identity's positions never exceeds its posted capital. Independent audits and pairwise-concavity checks are the remaining defense against collusion among identities that individually stay under any single reservation limit.
\end{pdSolution}
\begin{pdSolution}{bonded:ex:bonded-cartel-detection}
With $x = q_{\mathrm{floor}} - \mu$, the chapter's own closed form gives $\pi_C = x/3$, $\pi_D \to x$ as $\varepsilon/x\to0$, and $L=5x$, so
\[
p_d^\star=\frac{\pi_C-(1-\delta)\pi_D}{L+\delta\pi_D} = \frac{x/3-(1-\delta)x}{5x+\delta x} = \frac{1/3-(1-\delta)}{5+\delta}
\]
after cancelling the shared factor $x$. At $\delta=0.99$: $1/3-(1-0.99)=1/3-0.01=97/300$, and $5+0.99=5.99=599/100$, so $p_d^\star = (97/300)/(599/100) = 9700/179700 = 97/1797 \approx 0.05398$ [verified]---about $5.4\%$ per round, close to but distinct from the chapter's own worked point at $\delta=0.95$ ($p_d^\star\approx0.0478$): a more patient cartel ($\delta$ closer to $1$) needs a slightly \emph{higher} detection probability to stay fragile, because the discounted value of continuing the cartel is larger.

Cheap fresh identities shorten the effective horizon a punished member faces and let it shed the punishment entirely by re-registering---exactly the Sybil failure mode Exercise~\ref{bonded:ex:bonded-reputation-vs-bonds} names for reputation alone. A \emph{principal-bound} bond deters re-entry only when the expected forfeiture plus the onboarding and reputation-maturation cost of a fresh identity exceeds the gain from re-entering the cartel; a bond that is reusable across identities, or sized below that threshold, merely \emph{prices} re-entry rather than deterring it. The bond mechanism by itself never \emph{prevents} re-entry---nothing in this chapter's architecture stops a new registration---so the honest answer is: prices always, deters conditionally, prevents never.
\end{pdSolution}
\begin{pdSolution}{bonded:ex:bonded-insurance-alt-pricing}
This is a design obligation, not a closed theorem: no mechanism in this chapter proves a Pareto ordering for a posted-price catalog, so what follows is a defensible proposal, not a fictional derivation. A posted-price catalog quotes premiums by manifest risk class, principal history, exposure duration, and tail-capital band, recalibrated periodically against a hard reserve constraint---no per-transaction bidding, no auction latency. It conditionally improves on the static escrow of \S\ref{bonded:sec:scope-multiplier} exactly when the posted quote sits below the principal's idle-capital cost of self-insuring via a static bond, while the insurer remains individually rational and its reserves still cover the same tail risk the competitive market would have priced. Relative to \S\ref{bonded:sec:youle}'s bid-based mechanism, a posted catalog trades away instance-specific price discovery (Rothschild--Stiglitz's $q_I \to \mathbb{E}[d]$ convergence under competition) for lower manipulation surface and no per-transaction latency---the same trade the A6 cartel analysis makes salient in the other direction, since a catalog has no bid pattern for a colluding cartel to distort in the first place.

The proposal depends most on \emph{calibrated loss/tail dependence}: the catalog's bands must actually track the loss distribution within each class, and its capital reserve must be enforceably adequate against the tail, or the ``conditional'' improvement collapses into either under-pricing risk (insurer insolvency) or over-pricing it (pricing capable agents out, the same failure the threat-class bands of \S\ref{bonded:sec:pricing:threats} already warn against). The CC/NC/RP/NS assumptions of \S\ref{bonded:sec:youle:welfare} would need explicit class-conditional definitions before any formal Pareto claim could be stated for this alternative---this exercise does not supply that derivation, and neither should its answer key pretend to.
\end{pdSolution}
\begin{pdSolution}{bonded:ex:bonded-conservation-states}
The committed log (\path{proofs/bonded/conservation/Conservation.run.log}) reports:
\begin{lstlisting}
Model checking completed. No error has been found.
26818 states generated, 1716 distinct states found, 0 states left on queue.
The depth of the complete state graph search is 9.
\end{lstlisting}
These are two different numbers, and the difference is the exercise. TLC \emph{generates} a state every time it evaluates a transition, including transitions that land back on a state it has already seen by a different path---$26{,}818$ counts every such evaluation. The \emph{reachable} state count is the number of \textbf{distinct} states TLC actually found, $1{,}716$: every other generated state was a re-visit, not a new point in the state space. (Table~\ref{bonded:tab:verification-status} now carries both figures for this row; both are pinned from the same committed log, and only $1{,}716$ answers ``how many reachable states are there.'')

The invariant named in \texttt{Conservation.cfg}'s \texttt{INVARIANTS} block that gives the module its name is \texttt{Conservation}, defined as \texttt{TotalFree + TotalEscrow + burned = minted}---the free-form statement of Theorem~\ref{bonded:thm:conservation} at the granularity of individual mint/stake/slash/refund/payout operations rather than the three-bucket accounting state. The same run also checks \texttt{NoNegative} (no balance ever goes negative) and \texttt{Safety} (their conjunction); the log's ``No error has been found'' covers all three across every one of the $1{,}716$ reachable states, to search depth $9$.
\end{pdSolution}
\begin{pdSolution}{bonded:ex:bonded-conservation-weaken}
No run is needed to predict this: it follows from the definitions already in the spec. \texttt{Conservation} asserts $\mathit{TotalFree} + \mathit{TotalEscrow} + \mathit{burned} = \mathit{minted}$. Immediately after \texttt{Init} (all four quantities zero, satisfying the invariant trivially), the very first enabled \texttt{Mint}$(a, \mathit{amt})$ step increases $\mathit{TotalFree}$ by $\mathit{amt}$ while leaving $\mathit{minted}$ at $0$ under the stated weakening, so the left side becomes $\mathit{amt} > 0$ while the right side stays $0$. \texttt{Conservation} fails at depth~$1$---the shallowest possible counterexample, one step past the initial state, with no need to explore further.

\texttt{NoNegative}, by contrast, does \emph{not} catch this: every quantity involved (\texttt{free}, \texttt{escrow}, \texttt{burned}, \texttt{minted}) stays non-negative under the weakened \texttt{Mint}---the bug is a bookkeeping omission, not a sign violation, and \texttt{NoNegative}'s four non-negativity clauses are silent about the relationship \emph{between} the buckets. \texttt{Safety}, defined as \texttt{Conservation} $\wedge$ \texttt{NoNegative}, fails only because its \texttt{Conservation} conjunct does; this is exactly why the module checks \texttt{Conservation} as its own named invariant rather than folding it silently into \texttt{Safety} alone; a reader who only watched \texttt{NoNegative} would see nothing wrong.
\end{pdSolution}
\begin{pdSolution}{bonded:ex:bonded-magic-link-proverif}
Model a fresh token $t$, an issuance event $\mathrm{Issued}(u,t)$ tied to the user's mailbox, storage of an unguessable token (or a protected hash of it) with an expiry and an unused flag, delivery over a mailbox channel the attacker can read but not write into unprompted, redemption over an attacker-controlled network under \pdchapref{anchor}{The Anchor Protocol}'s Dolev-Yao assumptions, an atomic consume step, and recovery-key or session rotation on successful redemption.

The secrecy property is that the token (and any recovery secret it unlocks) stays unknown to the attacker absent mailbox compromise. The authentication property is injective: $\mathrm{Consumed}(u,t) \Rightarrow \mathrm{inj\text{-}event}(\mathrm{Issued}(u,t))$, exactly the form \S\ref{bonded:sec:federated-sovereign} already states for magic-link single-use---every consumption traces back to one, and only one, issuance, which is what rules out the race window between two concurrent redeemers. Add a concurrent-replay capability (two redemption attempts racing on the same token) and a mailbox-compromise capability (the attacker reads, but does not yet control, the delivery channel), neither of which the Phase~1 Harbor Card model needs: that model's attacker forges signatures or replays chains, but it has no notion of a \emph{delivery channel} an attacker might passively observe, because Harbor Cards are minted and verified entirely within the daemon's own boundary. Time and expiry also need an explicit abstraction here---ProVerif has no wall-clock semantics on its own---whereas the Harbor Card model's TTL is already folded into the token's signed payload.
\end{pdSolution}
\begin{pdSolution}{bonded:ex:bonded-gossip-proof}
This is an open mechanization obligation, not a solved one: Demers et al.\ is the epidemic-replication paper the chapter leans on, but no numbered theorem in it is stated in a form that yields this protocol's $\Theta(\log m)$ all-informed bound by direct substitution---the paper's push/pull rumor-spreading results assume a complete, uniform-peer overlay with reliable rounds, which is an assumption this chapter states but does not itself formalize. The mechanization obligation is therefore two steps, not one: first state the push/pull rumor-spreading theorem precisely, under exactly the complete-graph, independent-peer, reliable-round assumptions this chapter already uses (\S\ref{bonded:sec:claim-signaling-ic}'s ``observable history'' bullet leans on the same overlay); only then port that precise statement to a proof assistant. Citing the paper is not, by itself, a proof object.

On a general connected graph the $\Theta(\log m)$ bound does not transfer as stated: convergence time depends on the graph's conductance or mixing time, not on $m$ alone, and the honest replacement bound has a term roughly inverse in conductance, plus a tail probability rather than a clean deterministic count. A disconnected graph is the degenerate case that breaks the claim outright: gossip never converges across components that share no edge, no matter how long the protocol runs, because there is no path for a rumor to cross. A connected-but-bottlenecked mesh sits between the two: dissemination across the bottleneck cut is throttled by the cut's capacity, so a mechanized version of this bound would need to separate \emph{safety} (no revocation is ever falsely accepted or falsely dropped, which holds regardless of connectivity) from \emph{liveness} (every revocation eventually reaches every device, which needs the connectivity and mixing assumptions this section has so far only asserted in prose).
\end{pdSolution}
\begin{pdSolution}{bonded:ex:bonded-cuckoo-pollution}
An attacker with no special privilege submits a large volume of authenticated-looking but distinct revocation IDs---each one legitimate on its own, since the filter cannot distinguish a real revocation from a manufactured one by content alone. Each insertion fills buckets and, past a certain load, drives kick-chain failures in the filter's relocation step. Past the load at which Fan et al.'s false-positive bound holds, the filter is pushed outside the regime the bound was proved for: insertions the filter can no longer place are either rejected or force evictions of \emph{existing} entries, so a legitimate revocation already in the filter can be silently displaced.

The invariant that breaks is completeness of revocation---every previously revoked credential should still test positive---not the false-positive bound itself, which stays a probabilistic guarantee about \emph{non}-members. The $0.95$ ceiling defeats the attack described above only in the sense of preserving the regime the bound was proved in: refusing further insertions past that load keeps every already-inserted revocation intact, at the cost of refusing new ones (which is a capacity failure, loud and observable, rather than a silent integrity failure). It does not defeat a determined pollution attempt on its own, because an attacker willing to keep submitting IDs simply keeps the filter permanently at its ceiling, denying legitimate revocations a slot. Closing that requires what the ceiling alone does not provide: authenticated revocation provenance (only a party who can prove standing may insert), per-principal insertion quotas, an exact overflow log for anything the filter had to refuse, and a fail-closed policy---if insertion fails, the caller must not proceed as if the revocation succeeded. A probabilistic filter must never be allowed to silently drop an authoritative revocation.
\end{pdSolution}
